\ifdefined\included
\else
\documentclass[a4paper,11pt,twoside]{StyleThese}
\include{formatAndDefs}
\sloppy
\begin{document}
\fi


\chapter*{Introduction}
\addstarredchapter{Introduction} %Sinon cela n'apparait pas dans la table des matières

\section*{Contexte}

\section*{Problématique}

\section*{Approche Scientifique}

\section*{Contributions}

\section*{Contexte au sein du PEPR Réseaux du futur}
Dans le cadre de son programme de recherche scientifique France 2030, le gouvernement français a lancé plusieurs programmes PEPR (Programmes et Equipements Priotaires de Recherche) thématiques dans différents domaines scientifiques. L'un d'entre eux est consacré à la 5G et aux réseaux du futur (NF). Au sein de ce programme, NF-HiSec vise à traiter les questions de sécurité des réseaux du futurn, en particulier la 5G. Le projet NF-HiSec cherche à relever les principaux défis en matière de sécurité en se concentrant sur cinq objectifs principaux :

\begin{itemize}
\item \textbf{Détection et atténuation des attaques dans des environnements complexes :} Développer des techniques avancées pour identifier et atténuer les cyberattaques dans des réseaux dynamiques et en constante évolution.
\item \textbf{Protection des données personnelles, avec un accent sur l'interception légale :} Sécuriser les données personnelles même lors d'interceptions légales, tout en garantissant le respect des lois sur la vie privée.
\item \textbf{Modélisation des mécanismes de sécurité pour répondre aux besoins des applications :} Adapter les mesures de sécurité aux exigences spécifiques des différentes applications utilisant ces réseaux.
\item \textbf{Intégration de la sécurité à travers les frontières matériel et logiciel :} Formaliser la connexion entre les composants matériels et logiciels afin d'assurer la mise en œuvre de mesures de sécurité complètes à tous les niveaux.
\end{itemize}

Dans le cadre de cette large initiative, le projet NF-HiSec joue donc un rôle central dans le développement de méthodes et d'outils innovants visant à sécuriser de manière exhaustive les réseaux du futur. Cette thèse s'inscrit ainsi comme une contribution à un effort commun, cherchant à sécuriser une partie de l'ecosystème divers et complet qu'est la 5G.


\section*{Plan}



\ifdefined\included
\else
\bibliographystyle{acm}
\bibliography{These}
\end{document}
\fi